\subsection{Strategic adoption among incumbents}\label{sec:ext-game}

Proposition~\ref{prop:ext-timing} dates one firm's adoption against a falling cost, and Proposition~\ref{prop:ext-should}(a) compares outcomes once all rivals have made the same saving; but rivals choose their dates too. We model adoption as a game of timing with the canonical models of \citet{reinganum1981diffusion,reinganum1981market} and \citet{fudenberg1985preemption}, whose timing structure is that of Proposition~\ref{prop:ext-timing}, paired with the modern theory of pass-through \citep{weyl2013pass}, $n$-player preemption \citep{argenziano2014clustering} and the AI-native entrant of Assumption~\ref{ass:ext-entrant}. Proofs are in Appendix~\ref{app:proofs-game}; \texttt{code/verify\_part2\_competition.py} checks every claim and number. Examples are illustrative.

\begin{assumption}[Adoption game]\label{ass:ext-game}
There are $n_c\ge2$ incumbents. Each can adopt once, at a date of its choice, a bundle of automation that lowers its full unit cost by $\delta_c=\sum_j\mathcal B_jv_j/Y>0$ (Definition~\ref{def:ext-effcost}), where $Y$ is its output. Adopting at $T$ costs $F_A(T)=F_0e^{-\delta_FT}$, $F_0>0$, $\delta_F\ge0$, as in Proposition~\ref{prop:ext-timing}(a), and $F_A$ may include the other transition costs $\mathcal T$. Adoptions are observed at once. At date $t$ a firm in state $a\in\{0,1\}$ ($a=1$ once it has adopted), $k$ of whose rivals have adopted, earns $e^{Gt}\mathcal E_a(k)$, with $r>G>-\delta_F$. The private gain $\mathcal E^{\rm priv}(k):=\mathcal E_1(k)-\mathcal E_0(k)$ is positive and strictly decreasing in $k$ (adoptions are strategic substitutes), and $\mathcal E_1$ and $\mathcal E_0$ are strictly decreasing in $k$ (a rival's adoption hurts). The gain when all adopt is $\mathcal E^{\rm eq}:=\mathcal E_1(n_c-1)-\mathcal E_0(0)$, as in Proposition~\ref{prop:ext-should}(a). The single-firm date of a gain flow $\mathcal E$,
\begin{equation}\label{eq:ext-game-To}
T^\circ(\mathcal E):=\max\Bigl\{0,\ \frac1{G+\delta_F}\ln\frac{(r+\delta_F)F_0}{\mathcal E}\Bigr\}\ \ (\mathcal E>0),\qquad T^\circ(\mathcal E):=\infty\ \ (\mathcal E\le0),
\end{equation}
is the date at which $e^{Gt}\mathcal E$ first covers the user cost $(r+\delta_F)F_A(t)$ (Proposition~\ref{prop:ext-timing}(b)).
\end{assumption}

With two firms, let $T_2:=T^\circ(\mathcal E^{\rm priv}(1))$. Following \citet{fudenberg1985preemption}, write $V^{\rm L}(t)$ for the value at date 0 of a firm that adopts first at $t$ while its rival follows at its best date $\max\{t,T_2\}$, $V^{\rm F}(t)$ for the rival's value, and $V^{\rm J}(t)$ for each firm's value when both adopt at $t$. With $I(a,b):=\int_a^be^{-(r-G)s}ds$, for $t\le T_2$,
\begin{equation}\label{eq:ext-game-values}
\begin{aligned}
V^{\rm L}(t)&=\mathcal E_0(0)I(0,t)+\mathcal E_1(0)I(t,T_2)+\mathcal E_1(1)I(T_2,\infty)-e^{-rt}F_A(t),\\
V^{\rm F}(t)&=\mathcal E_0(0)I(0,t)+\mathcal E_0(1)I(t,T_2)+\mathcal E_1(1)I(T_2,\infty)-e^{-rT_2}F_A(T_2),\\
V^{\rm J}(t)&=\mathcal E_0(0)I(0,t)+\mathcal E_1(1)I(t,\infty)-e^{-rt}F_A(t),
\end{aligned}
\end{equation}
and $V^{\rm L}=V^{\rm F}=V^{\rm J}$ for $t\ge T_2$.

\begin{lemma}[Pass-through and the adoption incentives]\label{lem:ext-game-pass}
\begin{enumerate}
\item[(a)] (Linear demand.) Let $\bar\varphi^{\rm own}$ and $\bar\varphi^{\rm ind}$ be the level pass-through of a saving that no rival makes and of one that all make (Lemma~\ref{lem:ext-retained}(b)), and $\bar\varphi^\times:=(\bar\varphi^{\rm ind}-\bar\varphi^{\rm own})/(n_c-1)$ the effect of one rival's saving on a firm's price. If, with unit margin mg before adoption, $\mathcal E_1(k)=[\mathrm{mg}+(1-\bar\varphi^{\rm own})\delta_c-k\bar\varphi^\times\delta_c]^2/B_q$ and $\mathcal E_0(k)=[\mathrm{mg}-k\bar\varphi^\times\delta_c]^2/B_q$, with $\bar\varphi^\times>0$ and $\mathrm{mg}>(n_c-1)\bar\varphi^\times\delta_c$, then the conditions of Assumption~\ref{ass:ext-game} on $\mathcal E_a$ hold. The Cournot market of Lemma~\ref{lem:ext-retained} has this form with $B_q=B$, $\bar\varphi^{\rm own}=\bar\varphi^\times=1/(n_c+1)$ and $\bar\varphi^{\rm ind}=n_c/(n_c+1)$; so does differentiated Bertrand competition with demand $q_i=A_B-p_i+s_B\sum_{l\ne i}p_l$, $0<s_B<1/(n_c-1)$, with $B_q=1$, $\bar\varphi^{\rm ind}=1/[2-(n_c-1)s_B]$ and $\bar\varphi^\times=s_B\bar\varphi^{\rm ind}/(2+s_B)$.
\item[(b)] (Cournot with general demand.) Let inverse demand $P(Q)$ be twice continuously differentiable with $P'<0$, all firms have unit cost $\hat c$ and output $Y$, $\mathsf e:=-QP''(Q)/P'(Q)$ be the curvature of inverse demand, and the symmetric equilibrium be regular, $n_c+1-\mathsf e>0$. Then $\bar\varphi^{\rm ind}=n_c/(n_c+1-\mathsf e)$, $\bar\varphi^{\rm own}=\bar\varphi^{\rm ind}/n_c$, Lemma~\ref{lem:prices}'s elasticity is $\varphi=\bar\varphi^{\rm ind}\hat c/P$, and as $\delta_c\to0$, for every $k$,
\begin{equation}\label{eq:ext-game-first}
\mathcal E^{\rm eq}=Y\delta_c\Bigl[1-\bar\varphi^{\rm ind}\Bigl(1-\frac1{n_c}\Bigr)\Bigr]+o(\delta_c),\qquad \mathcal E^{\rm priv}(k)=Y\delta_c\Bigl[\frac{2n_c-1}{n_c}-\frac{(n_c-1)\bar\varphi^{\rm ind}}{n_c^2}\Bigr]+o(\delta_c).
\end{equation}
The first is the incidence formula of \citet{weyl2013pass} with Cournot conduct $1/n_c$.
\item[(c)] (Constant elasticity.) If $Q=DP^{-\epsilon}$ with $n_c\epsilon>1$, then $P=n_c\epsilon\hat c/(n_c\epsilon-1)$, $\bar\varphi^{\rm ind}=n_c\epsilon/(n_c\epsilon-1)$, $\varphi=1$ and the Lerner index is $1/(n_c\epsilon)$. Earnings and revenue are proportional to $\hat c^{\,1-\epsilon}$, so $\mathcal E^{\rm eq}=\mathcal E_0(0)[(1-\delta_c/\hat c)^{1-\epsilon}-1]$ has the sign of $\epsilon-1$. The private gain in \eqref{eq:ext-game-first} is positive if and only if $\epsilon>(2n_c-1)/(2n_c^2-2n_c+1)$.
\end{enumerate}
\end{lemma}

\noindent\emph{Intuition and example.} Pass-through decides how much of a saving leaves the industry; the volume response decides how much a firm gains by taking business from rivals. With four Cournot firms the two pass-throughs are 20\% and 80\%, and $\varphi=0.64$ at a price of 1.00 and unit cost 0.80; under Bertrand competition with $s_B=0.2$ they are 0.519 and 0.714. Under constant elasticity $\epsilon=0.8$ with four firms, $\bar\varphi^{\rm ind}=1.455$: a common saving of 0.10 on a unit cost of 0.6875 cuts the price by 14.5\% and every firm's earnings by 3.1\%, though adopting alone still pays (the threshold in (c) is 0.28 with four firms, 0.6 with two).

\begin{proposition}[Diffusion and preemption]\label{prop:ext-game-preempt}
Under Assumption~\ref{ass:ext-game}:
\begin{enumerate}
\item[(a)] (Precommitment.) Let $T_k:=T^\circ(\mathcal E^{\rm priv}(k-1))$, $k=1,\dots,n_c$. If firms commit to dates at date 0, the dates $T_1\le\dots\le T_{n_c}$, assigned in any order, form a Nash equilibrium. They increase strictly where positive, so identical firms adopt in sequence \citep{reinganum1981diffusion,reinganum1981market}. The $k$-th adopter moves before $T^\circ(\delta_cY)$, the date of a firm that takes prices and volume as given, if and only if $\mathcal E^{\rm priv}(k-1)>\delta_cY$.
\item[(b)] (Preemption.) Let $n_c=2$, $T_2>0$ and $V^{\rm L}(0)<V^{\rm F}(0)$, and let $T_1:=T^\circ(\mathcal E^{\rm priv}(0))$ and $T_D:=T^\circ(\mathcal E_1(0)-\mathcal E_0(1))$. Then $0<T_D<T_1<T_2$, and $V^{\rm L}-V^{\rm F}$ rises strictly on $[0,T_D]$, falls strictly on $[T_D,T_2]$ and vanishes at $T_2$, so it has a unique root $T^{\rm P}\in(0,T_D)$. If firms can respond to each other, the timing game of \citet{fudenberg1985preemption} has a subgame-perfect equilibrium in which one firm adopts at $T^{\rm P}$ and the other at $T_2$. Rents are equalized at $V^{\rm L}(T^{\rm P})=V^{\rm F}(T^{\rm P})$, and the race costs the leader $V^{\rm L}(T_1)-V^{\rm L}(T^{\rm P})>0$.
\end{enumerate}
\end{proposition}

\noindent\emph{Intuition and example.} A committed firm adopts when its private gain covers the user cost, and since each adoption lowers the next firm's gain, the dates fan out. A firm that can be preempted cannot wait for its best date: while leading is worth more than following, a rival would move first, so the first adoption is pulled forward until the two are worth the same. Take four Cournot incumbents, each selling 20 million units a year at \$1.00 with unit cost \$0.80 and able to save \$0.10 a unit, with $r=0.10$, $G=0.04$ and, as in Proposition~\ref{prop:ext-timing}'s example, an adoption cost of \$20 million falling 30\% a year. The private gains are \$3.84, 3.52, 3.20 and 2.88 million a year (Proposition~\ref{prop:ext-should}(a) with $B=0.01$), and the precommitment dates 2.16, 2.41, 2.69 and 3.00 years, all before the 4.08 years at which a firm ignoring its rivals would adopt the \$2.0 million engineering saving: the private gain includes business taken from rivals. The durable gain of \$0.84 million would cover the user cost only after 6.63 years. In a duopoly with the same numbers per firm, $(\mathcal E_0(0),\mathcal E_1(0),\mathcal E_0(1),\mathcal E_1(1))=(4.00,7.11,2.78,5.44)$, $T_D=1.80$, $T_1=2.78$ and $T_2=3.23$, and preemption brings the first adoption forward to $T^{\rm P}=0.63$ years; each firm is worth \$78.2 million, against \$81.1 million for an unchallenged leader at $T_1$ and \$66.7 million if no one adopted. For $n_c>2$ we report the precommitment dates, because rent equalization does not extend to oligopoly \citep{fudenberg1985preemption} and $n$-player preemption can cluster investments \citep{argenziano2014clustering}.

\begin{proposition}[Waiting games]\label{prop:ext-game-wait}
Under Assumption~\ref{ass:ext-game} with $n_c=2$ and $T_2>0$:
\begin{enumerate}
\item[(a)] (Late joint adoption.) Let $T_J:=T^\circ(\mathcal E^{\rm eq})$, which exceeds $T_2$ and is infinite if $\mathcal E^{\rm eq}\le0$. For finite $T\in[T_2,T_J]$, joint adoption at $T$ is a subgame-perfect equilibrium if and only if $V^{\rm J}(T)\ge V^{\rm L}(T_1)$, and no joint date outside $[T_2,T_J]$ is. Such equilibria exist (never adopting, if $T_J=\infty$) if and only if
\begin{equation}\label{eq:ext-game-joint}
\int_{T_2}^{T_J}e^{-rt}\bigl[(r+\delta_F)F_A(t)-\mathcal E^{\rm eq}e^{Gt}\bigr]dt\ \ge\ \int_{T_1}^{T_2}e^{-rt}\bigl[\mathcal E^{\rm priv}(0)e^{Gt}-(r+\delta_F)F_A(t)\bigr]dt,
\end{equation}
and they give each firm more than the preemption equilibrium.
\item[(b)] (Learning.) Let the bundle work for the industry with known probability $\varpi\in(0,1)$, the first adoption reveal at once whether it does, and a failure leave the adopter's cost unchanged, its outlay sunk and no further adoption. Then $V^{\rm L}_\varpi(t):=\varpi V^{\rm L}(t)+(1-\varpi)[\mathcal E_0(0)/(r-G)-e^{-rt}F_A(t)]$ and $V^{\rm F}_\varpi(t):=\varpi V^{\rm F}(t)+(1-\varpi)\mathcal E_0(0)/(r-G)$. If $\varpi\le\bar\varpi:=\mathcal E^{\rm priv}(1)/[\mathcal E_1(0)-\mathcal E_0(1)]\in(0,1)$, then $V^{\rm F}_\varpi>V^{\rm L}_\varpi$ at every date: no firm ever gains by preempting.
\item[(c)] (War of attrition.) Under (b), let $\varpi\le\bar\varpi$, $\mathcal E^{\rm eq}>0$ and $T^{\rm W}:=T^\circ(\varpi\mathcal E^{\rm eq})\ge T_2$. Then $V^{\rm L}_\varpi$ rises strictly before $T^{\rm W}$ and falls strictly after it. (i) For either assignment of roles there is an equilibrium in which the leader adopts at $T^{\rm W}$ and the follower adopts at once if the bundle works and never otherwise. (ii) There is a symmetric equilibrium in which, until someone adopts, each firm adopts at the hazard rate
\begin{equation}\label{eq:ext-game-hazard}
\nu^{\rm W}(t)=\frac{\varpi\mathcal E^{\rm eq}e^{Gt}-(r+\delta_F)F_A(t)}{(1-\varpi)F_A(t)}\quad(t\ge T^{\rm W}),\qquad \nu^{\rm W}(t)=0\quad(t<T^{\rm W});
\end{equation}
each firm is worth $V^{\rm L}_\varpi(T^{\rm W})$, less than the follower in (i), and the first adoption comes a random delay $s$ after $T^{\rm W}$, with $\Pr(s>u)=\exp\{-2\int_{T^{\rm W}}^{T^{\rm W}+u}\nu^{\rm W}\}$.
\end{enumerate}
\end{proposition}

\noindent\emph{Intuition and example.} Two forces make firms wait. If the gain from leading is small, both can wait for the late date at which the durable gain covers the user cost, since a firm that broke ranks would gain less from leading than it loses by triggering early adoption \citep{fudenberg1985preemption}. If success is uncertain, the first adopter reveals the answer to its rival, who copies a success and avoids a failure; following beats leading, and adoption becomes a war of attrition \citep{hendricks1988war} with second-mover advantages \citep{hoppe2000second}, timed by the risk-adjusted durable gain $\varpi\mathcal E^{\rm eq}$. In the duopoly above, joint adoption at any date in $[3.32,5.03]$ is an equilibrium worth up to \$81.8 million to each firm, against \$78.2 million under preemption. Such equilibria exist here at every cut for which the follower waits ($T_2>0$); they fail for large cuts and costs, and with a cut of \$0.31 and an adoption cost of \$80 million the outcome is preemption at 0.66 years. If the bundle works with probability $\varpi=0.5$, below $\bar\varpi=0.615$ (preemption reappears only above 0.893), the war of attrition starts at $T^{\rm W}=7.07$ years, and the first adoption follows after a mean delay of 1.52 years (median 1.46); each firm is worth \$73.4 million, against \$74.0 million for an assigned follower.

\begin{proposition}[The automation prisoner's dilemma]\label{prop:ext-game-pd}
\begin{enumerate}
\item[(a)] (Static.) Let $n_c$ symmetric firms each decide whether to adopt at an annualized cost $\mathcal K$, with $\mathcal E^{\rm priv}(k)$ decreasing and $\mathcal E_0(k)$ strictly decreasing in $k$. Adopting is strictly dominant if and only if $\mathcal K<\mathcal E^{\rm priv}(n_c-1)$, and when all adopt each firm is no better off if and only if $\mathcal K\ge\mathcal E^{\rm eq}$. The dilemma window $[\mathcal E^{\rm eq},\mathcal E^{\rm priv}(n_c-1))$ has width $\mathcal E_0(0)-\mathcal E_0(n_c-1)>0$; to first order in symmetric Cournot it is
\begin{equation}\label{eq:ext-game-window}
\frac{\mathcal K}{Y\delta_c}\in\Bigl[1-\bar\varphi^{\rm ind}\Bigl(1-\frac1{n_c}\Bigr),\ \frac{2n_c-1}{n_c}-\frac{(n_c-1)\bar\varphi^{\rm ind}}{n_c^2}\Bigr),
\end{equation}
which under linear demand is $[2/(n_c+1),2n_c/(n_c+1))$: only the share $1/n_c$ of the private gain is durable.
\item[(b)] (Pass-through, elasticity and the number of firms.) To first order, common adoption lowers every firm's gross earnings, so that the dilemma holds at every $\mathcal K\ge0$ at which adopting is privately worthwhile, if and only if $\bar\varphi^{\rm ind}>n_c/(n_c-1)$, that is $\mathsf e>2$, whatever $n_c$. Under constant elasticity it holds exactly if and only if $\epsilon<1$; under linear demand, never.
\item[(c)] (Dynamic.) Under Assumption~\ref{ass:ext-game}, in the precommitment equilibrium the last adopter is no better off than if no firm ever adopted whenever $\mathcal E^{\rm eq}\le(r-G)e^{-GT_{n_c}}F_A(T_{n_c})$, which for $T_{n_c}>0$ reads
\begin{equation}\label{eq:ext-game-dyn}
\mathcal E^{\rm eq}\le\frac{r-G}{r+\delta_F}\,\mathcal E^{\rm priv}(n_c-1).
\end{equation}
Under the conditions of Proposition~\ref{prop:ext-game-preempt}(b), the same condition with $n_c=2$ leaves both duopolists no better off in the preemption equilibrium. Under linear demand it holds, to first order, when $n_c\ge(r+\delta_F)/(r-G)$.
\item[(d)] (Rule of $X$.) In symmetric Cournot with $b=0$, let $S_0$, $S_1$ be a firm's revenue before and after common adoption and $m_0$ its margin before. Its margin net of $\mathcal K$ changes by exactly $[\mathcal E^{\rm eq}-\mathcal K-m_0(S_1-S_0)]/S_1$. To first order, revenue changes by $-(1-\epsilon)\varphi\delta_c/\hat c$ (recorded once in growth, with weight $X$) and the margin before $\mathcal K$ by $(1-\varphi)\delta_c/P$, and a continuing fall in unit cost at rate $\delta_u$ lowers nominal growth by $(1-\epsilon)\varphi\delta_u$ (Lemma~\ref{lem:prices}). Under constant elasticity the margin falls by exactly $\mathcal K/S_1$; if also $\epsilon<1$, every incumbent's Rule-of-$X$ score falls in both terms, and its value $Mv(s)$ (Proposition~\ref{prop:dynamic}(a)) falls twice over at unchanged shares, as $M$ shrinks by the factor $(1-\delta_c/\hat c)^{1-\epsilon}$ and $v$ falls with the margin intercept.
\end{enumerate}
\end{proposition}

\noindent\emph{Intuition and example.} Each firm adopts because a saving it alone made would take business from its rivals. When all adopt, no business moves and the price falls, leaving the durable share: small when pass-through is high and firms are many, negative when pass-through exceeds $n_c/(n_c-1)$, because industry earnings then fall with industry costs \citep{weyl2013pass}. In the four-firm example the window is $[0.84,2.88)$ million dollars a year, against $[0.8,3.2)$ to first order; under constant elasticity $\epsilon=0.8$, common adoption costs every firm 3.1\% of earnings even if free. The dynamic test compares the durable gain with the share $(r-G)/(r+\delta_F)=0.15$ of the last adopter's private gain, the annualized cost of its investment at its chosen date. With four firms the ratio is 0.29 and the last adopter gains; with ten it is 0.125 and the last adopter loses \$2.73 million (the first-order bound is $n_c\ge6.67$). Falling adoption costs soften the dilemma: with a constant cost of \$20 million all four firms adopt at once, and each loses \$6.0 million although the last adopter's private gain is worth \$48 million. On the Rule of 40, common adoption in the linear market raises revenue by 1.2\% and the margin before adoption cost from 20.0\% to 23.9\%, but at $\mathcal K=\mathcal E^{\rm eq}$ the margin is 19.8\%; under constant elasticity $\epsilon=0.8$, revenue falls by 3.1\% with the margin of 31.25\% unchanged, and a continuing 5\% annual fall in unit cost costs 1 point of growth. Faster growth of AI-investing firms \citep{babina2024artificial} fits the business-stealing part of the private gain, which is not durable for the industry.

\begin{proposition}[AI-native entry and defensive automation]\label{prop:ext-game-entry}
Let the incumbents be the Cournot firms of Lemma~\ref{lem:ext-game-pass}(a), with price $P_0$ and unit cost $\hat c_0$ before AI and $P_k=P_0-k\delta_c/(n_c+1)$ once $k$ have adopted, and the entrants as in Assumption~\ref{ass:ext-entrant}, with full unit cost $\hat c_E(t)$ strictly decreasing to $\hat c_E^\infty$; write $P_E(t):=\hat c_E(t)+p_{\rm sw}$ and $P_E^\infty:=\hat c_E^\infty+p_{\rm sw}$. Entry binds when the incumbents' price exceeds $P_E(t)$ (Proposition~\ref{prop:ext-should}(b)); each incumbent then sells its installed base $Ye^{Gt}$ at $P_E(t)$.
\begin{enumerate}
\item[(a)] (Entry ends the game.) From $\bar T_M:=\inf\{t:P_E(t)<P_{n_c}\}$ on, entry binds whatever incumbents do, and an incumbent's gain from adopting is $\delta_cYe^{Gt}$ whatever its rivals do. Pass-through is zero, the motives of Propositions~\ref{prop:ext-game-preempt} and \ref{prop:ext-game-wait}(a) disappear (the learning externality of Proposition~\ref{prop:ext-game-wait}(b) remains), and an incumbent that stays and has not adopted at $t\ge\bar T_M$ adopts at $\max\{t,T^\circ(\delta_cY)\}$.
\item[(b)] (Deterrence is a public good.) The breakpoint $T_M(k):=\inf\{t:P_E(t)<P_k\}$ is nondecreasing in the number of adopters $k$, so each adoption postpones entry for all incumbents. If $P_{n_c}\le P_E^\infty$, entry never binds if and only if, for each $j\le\bar k:=\min\{k:P_k\le P_E^\infty\}$, the $j$-th adoption takes place by $T_M(j-1)$.
\item[(c)] (Defensive automation.) Let $P_E^\infty<P_{n_c}$ and let adoption cost $\mathcal K$ a year. In the long run an adopter's flow exceeds a non-adopter's by $\delta_cY-\mathcal K$, since it keeps all of its saving, and exceeds its own flow before AI by $(P_E^\infty-P_0+\delta_c)Y-\mathcal K$. Adoption is therefore defensive, in the sense of Proposition~\ref{prop:ext-should}(a), if and only if $(P_E^\infty-P_0+\delta_c)Y\le\mathcal K<\delta_cY$, and if $P_E^\infty\le P_0-\delta_c$ every adoption worth making is defensive. It is needed for viability if and only if $P_E^\infty<\hat c_0$, and it secures viability if and only if $P_E^\infty\ge\hat c_0-\delta_c+\mathcal K/Y$ (Proposition~\ref{prop:ext-should}(c)).
\item[(d)] (Free replication.) If entrants make the incumbent's saving at least as cheaply, $\hat c_E^\infty\le\hat c_0-\delta_c$, as under the free replication of Proposition~\ref{prop:entry}(b), the incumbent's long-run unit margin is at most $p_{\rm sw}$ plus any cost advantage from a scarce input (Proposition~\ref{prop:entry}(c)).
\end{enumerate}
\end{proposition}

\noindent\emph{Intuition and example.} Once entry binds, entrants' cost sets the price: an incumbent keeps its whole saving but takes no business from rivals, so the strategic motives vanish. The entrant also moves the benchmark: relative to the world before AI, automation no longer raises the incumbent's margin; it only limits the fall. Let entrants' cost be $\hat c_E(t)=0.70+0.50e^{-0.4t}$, with $p_{\rm sw}=0.05$, against the four incumbents. Entry binds at 1.73, 1.94, 2.17, 2.42 or 2.70 years if zero to four incumbents have adopted, and from 2.70 years whatever they do; no schedule deters it, since $P_4=0.92>P_E^\infty=0.75$. After 2.70 years the gain from adopting is the \$2.0 million engineering saving, so a firm that has not adopted waits until 4.08 years, before its margin turns negative at 5.76 years. In the long run an adopter earns a unit margin of 0.05, against 0.20 before AI and $-0.05$ without automation; since $P_E^\infty=0.75\le P_0-\delta_c=0.90$, every adoption worth making is defensive (window $[-3.0,2.0)$ million dollars a year). Proposition~\ref{prop:prod-cannibal}(c) gives the analogous product choice.

\begin{corollary}[Durable gain, concentration and should versus could]\label{cor:ext-game-engine}
Let HHI be the incumbents' Herfindahl index, with $n_c=1/\mathrm{HHI}$ for symmetric firms, under the conditions of Lemma~\ref{lem:ext-game-pass}(b).
\begin{enumerate}
\item[(a)] (Durable share.) To first order a firm keeps the share
\begin{equation}\label{eq:ext-game-dur}
\mathrm{dur}:=1-\bar\varphi^{\rm ind}(1-\mathrm{HHI})=\frac{2-\mathsf e}{1/\mathrm{HHI}+1-\mathsf e}
\end{equation}
of a saving $\delta_cY$ that every rival makes, measured in dollars of profit with its volume gain counted (it is not a margin share: at the lower price the margin keeps $1-\bar\varphi^{\rm ind}u/p$ of the saving to first order, which does not depend on HHI; Proposition~\ref{prop:tasks-aggregate}(a)), and $\mathrm{priv}:=2-\mathrm{HHI}-\mathrm{HHI}(1-\mathrm{HHI})\bar\varphi^{\rm ind}$ of one that none makes ($2\,\mathrm{HHI}/(1+\mathrm{HHI})$ and $2/(1+\mathrm{HHI})$ under linear demand; $\mathrm{dur}=\mathrm{HHI}(\epsilon-1)/(\epsilon-\mathrm{HHI})$ under constant elasticity), and $(1-\mathrm{rep})\,\mathrm{priv}+\mathrm{rep}\cdot\mathrm{dur}$ if the share rep of rivals can replicate it. With constant curvature, dur is one under monopoly and rises with concentration if and only if $\mathsf e<2$. Where entry binds (Proposition~\ref{prop:ext-game-entry}(c)), $\mathrm{dur}=(P_E^\infty-P_0+\delta_c)/\delta_c$ and $\mathrm{priv}=1$.
\item[(b)] (Classification.) For a bundle with advantage $\sum_j\mathcal B_jv_j=\delta_cY$, success probability $\varpi$ and transition cost $\mathcal T$, let $\mathrm{NPV}^{\rm dur}:=\varpi\,\mathrm{dur}\,\delta_cY/(r-G)-\mathcal T$ and $\mathrm{NPV}^{\rm priv}:=\varpi\,\mathrm{priv}\,\delta_cY/(r-G)-\mathcal T$. The bundle is value-creating if $\mathrm{NPV}^{\rm dur}>0$; defensive if $\mathrm{NPV}^{\rm dur}\le0<\mathrm{NPV}^{\rm priv}$ (the strategic case of Proposition~\ref{prop:ext-should}(a)) or if it is needed for viability (Definition~\ref{def:ext-motives}(ii)); and feasible but not worth doing otherwise. Its date is the strategic date of Propositions~\ref{prop:ext-game-preempt} and \ref{prop:ext-game-wait}, capped by the deadline of Proposition~\ref{prop:ext-should}(e).
\item[(c)] (Expected durable gain.) The expected durable gain is $\varpi\,\mathrm{dur}\,\delta_cY-\mathcal K$ a year; when $\mathsf e<2$ it rises with concentration, from $-\mathcal K$ in a fragmented industry to $\varpi\delta_cY-\mathcal K$ under monopoly.
\end{enumerate}
\end{corollary}

\begin{table}[t]
\centering
\caption{Durable and private shares of a replicable saving, to first order (Corollary~\ref{cor:ext-game-engine}(a))}\label{tab:ext-game-dur}
\footnotesize
\begin{tabular}{@{}lcccccc@{}}
\toprule
Number of incumbents $n_c$ (HHI $=1/n_c$) & 1 & 2 & 3 & 4 & 6 & 10\\
\midrule
Linear demand: durable share $2/(n_c+1)$ & 1.00 & 0.67 & 0.50 & 0.40 & 0.29 & 0.18\\
Linear demand: private share $2n_c/(n_c+1)$ & 1.00 & 1.33 & 1.50 & 1.60 & 1.71 & 1.82\\
Constant elasticity $\epsilon=2$: durable share & 1.00 & 0.33 & 0.20 & 0.14 & 0.09 & 0.05\\
Constant elasticity $\epsilon=0.8$: durable share & n.a. & $-$0.33 & $-$0.14 & $-$0.09 & $-$0.05 & $-$0.03\\
\bottomrule
\end{tabular}

\smallskip
\begin{minipage}{0.92\linewidth}\footnotesize
Notes: shares of the engineering saving $\delta_cY$. A monopoly with $\epsilon<1$ has no finite price (n.a.). Values computed by \texttt{code/verify\_part2\_competition.py}.
\end{minipage}
\end{table}

\noindent\emph{Intuition and example.} The engine's question is not whether a saving can be made but how much of it survives the rivals' response. A monopolist keeps all of it; four Cournot firms with linear demand keep 40\% once all adopt, though each would gain 160\% by adopting alone; under inelastic demand the share is negative (Table~\ref{tab:ext-game-dur}). A bundle that saves the example's \$2.0 million a year, works with probability $\varpi=0.6$ and costs $\mathcal T=\$20$ million has $\mathrm{NPV}^{\rm dur}=-\$12.0$ million and $\mathrm{NPV}^{\rm priv}=\$12.0$ million: it is defensive, to be done only because rivals will do it. Definition~\ref{def:ext-motives}, which counts the retained share $1/(n_c+1)$ of the unit margin at fixed volume, gives $-\$16.0$ million; the volume response doubles the durable share under linear demand. For should versus could, three implications follow. In fragmented markets most replicable automation is defensive, worth doing to avoid a loss rather than to raise the margin. In concentrated markets more of it is durable, but duopolies may delay it, by tacit joint waiting or, when success is uncertain, by a war of attrition. Where AI-native entry binds, the durable gain is set by the entrants' cost and the switching premium, and the incumbents' game no longer matters. The results assume symmetric incumbents, one bundle per firm and deterministic adoption costs; with asymmetric firms the engine uses $n_c=1/\mathrm{HHI}$ as an approximation.

\noindent\emph{Calibration of the industry pass-through.} Lemma~\ref{lem:ext-game-pass}(b) derives $\bar\varphi^{\rm ind}=n_c/(n_c+1-\mathsf e)$ for homogeneous Cournot firms, which is 0.8 for four firms under linear demand and at least 0.75 for three or more. That is a benchmark of the model, not a measurement, and it lies above the evidence on cost pass-through. Developers of differentiated apps cut prices on 9\% of products after a cost cut of about ten points of revenue, with a median cut of 2.5\%, within three months (a study funded by Apple, without a control group) \citep{choi2025dma}; about 70\% of industry-wide energy cost changes pass through to U.S.\ manufacturing prices in the short to medium run \citep{ganapati2020energy}; and final pass-through of value-added tax is about 100\% for standard-rate changes and about 30\% for reduced-rate changes in the euro area \citep{benedek2015vat}. Part~II therefore treats $\bar\varphi^{\rm ind}$ as a scenario dimension with central value 0.5, the medium-run case between the software and manufacturing anchors, and scenarios 0 (fixed prices, the short run for differentiated software) and 0.7 (an industry-wide cost shock); values near one belong to undifferentiated services in the long run (all class b). Equation~\eqref{eq:ext-game-dur} is kept with the measured value in place of the derived one; within the lemma's model, 0.5 and 0.7 at four firms correspond to curvatures $\mathsf e=n_c+1-n_c/\bar\varphi^{\rm ind}$ of $-3$ and $-0.71$, that is, to more concave demand, and they stand for the differentiation and frictions the homogeneous model leaves out. At four symmetric firms ($\mathrm{HHI}=0.25$) the durable share is 1.0, 0.625 and 0.475 at pass-through 0, 0.5 and 0.7, against 0.40 at the Cournot benchmark, and the private share is 1.75, 1.66 and 1.62 against 1.60. The bundle of the example then has $\mathrm{NPV}^{\rm dur}$ of \$0, $-\$7.5$ million and $-\$10.5$ million and $\mathrm{NPV}^{\rm priv}$ of \$15.0, \$13.1 and \$12.4 million: it stays defensive at every scenario, and at fixed prices it is on the boundary, worth exactly its cost. These are dollar shares. The rate $\bar\varphi^{\rm ind}$ also sets the margin, and so the score: at the lower price a replicated saving keeps $1-\bar\varphi^{\rm ind}u/p$ of its margin gain to first order, whatever the concentration (Part~I's mapping, which applies the rate to every unit-cost saving and computes the margin at the post-cut price), so a lower pass-through raises the rank of cost-of-revenue automation among the levers of Proposition~\ref{prop:playbook-rank} through that share, not through the durable share; the tables above remain the linear-demand benchmark.
